\chapter{Planificación temporal}
\label{chap:planificacion-temporal}

\begingroup
\small
\setstretch{1.05}
\setlength{\parskip}{0.3\baselineskip}
\newcolumntype{Y}{>{\raggedright\arraybackslash}X}
\newcolumntype{C}[1]{>{\centering\arraybackslash}p{#1}}
\newcolumntype{L}[1]{>{\raggedright\arraybackslash}p{#1}}
\definecolor{upcblue}{HTML}{0076A8}
\definecolor{deepblue}{HTML}{174A6E}
\definecolor{green}{HTML}{5D8A66}
\definecolor{amber}{HTML}{D98E04}

La planificación cubre el proyecto completo, desde la definición del protocolo
experimental hasta la defensa. Se adopta un enfoque iterativo e incremental:
primero se obtiene un producto mínimo viable (MVP) con datos, línea base y
evaluación; después se incorporan tres familias visuales; finalmente se comparan
precisión y coste, se selecciona un modelo y se valida su inferencia en un
entorno limitado. La documentación y el seguimiento se realizan en paralelo.

\begin{table}[H]
  \centering
  \begin{tabularx}{\textwidth}{>{\bfseries}L{4.2cm}Y}
    \toprule
    Inicio planificado & 14/09/2026 \\
    Fin técnico del proyecto & 15/01/2027 \\
    Fin desde el punto de vista de la FIB & 29/01/2027 (objetivo provisional) \\
    Lectura prevista del TFG & 29/01/2027, pendiente de confirmación en Racó (TBD) \\
    Esfuerzo total & 540 horas, equivalentes a 90 jornadas de 6 horas \\
    Ritmo ordinario & 6 h/día, de lunes a viernes; 30 h/semana \\
    Horizonte temporal & 20 semanas naturales; la diferencia entre capacidad y
    esfuerzo absorbe festivos, indisponibilidad y riesgo \\
    Responsable & Estudiante, con supervisión del director y codirector \\
    \bottomrule
  \end{tabularx}
  \caption[Datos globales de la planificación]{Datos globales de la
  planificación temporal. Elaboración propia.}
  \label{tab:planificacion-datos-globales}
\end{table}

Las horas incluyen implementación, comprobación, registro y cierre de cada
tarea. Los entrenamientos se estiman por tiempo de trabajo activo: el cómputo
desatendido no cuenta como dedicación, aunque condiciona el calendario. Ninguna
tarea técnica alcanza 50 horas. La memoria se mantiene como una tarea
transversal de 60 horas porque su actividad es homogénea y recurrente.

El trabajo se divide en iteraciones de dos semanas. Al final de cada una se
revisan los resultados, los riesgos y el trabajo pendiente en el tablero Kanban.
El Gantt fija el marco global y el tablero permite reordenar tareas sin perder
los hitos ni la fecha final.

Los recursos se codifican así: R1, estudiante; R2, director y codirector; R3,
portátil, Python, TensorFlow/Keras y Git; R4, Nutrition5k, almacenamiento y
entorno de cómputo; R5, LiteRT y dispositivo \emph{edge} concreto (TBD); y R6,
Notion, LaTeX y herramientas FIB.

\clearpage
\section{Tareas, estimaciones y recursos}

Las tablas~\ref{tab:planificacion-tareas-1} y
\ref{tab:planificacion-tareas-2} describen tareas verificables, no simples
fases. El resultado de cierre concreta qué debe existir para considerar
terminada cada actividad; ``Pred.'' identifica sus predecesoras.

\begin{table}[H]
  \centering
  \scriptsize
  \begin{tabularx}{\textwidth}{C{0.75cm}Y C{0.55cm} L{1.35cm} L{1.75cm}}
    \toprule
    \textbf{Cód.} & \textbf{Tarea y resultado de cierre} & \textbf{h} &
    \textbf{Pred.} & \textbf{Recursos} \\
    \midrule
    T01 & \textbf{Fijar alcance y protocolo.} Objetivos, variables, métricas y criterios de comparación aprobados. & 14 & -- & R1, R2, R6 \\
    T02 & \textbf{Revisar bibliografía dirigida.} Matriz de trabajos, modelos, datos, métricas y límites relevantes. & 24 & -- & R1, R3, R6 \\
    T03 & \textbf{Preparar el entorno reproducible.} Repositorios, dependencias, pruebas y estructura experimental operativos. & 10 & T01 & R1, R3 \\
    T04 & \textbf{Adquirir y validar Nutrition5k.} Descarga controlada, inventario, integridad y muestra inspeccionada. & 18 & T01 & R1, R3, R4 \\
    T05 & \textbf{Construir el flujo de datos.} Extracción de fotogramas, carga, transformación y caché reproducibles. & 28 & T03, T04 & R1, R3, R4 \\
    T06 & \textbf{Definir particiones y controles de fuga.} División fijada por plato y comprobaciones automatizadas. & 14 & T05 & R1, R3, R4 \\
    T07 & \textbf{Implementar y evaluar la línea base.} Predictor medio, métricas y artefactos de referencia guardados. & 14 & T06 & R1, R3, R4 \\
    T08 & \textbf{Implementar la CNN de regresión.} Arquitectura y cinco salidas integradas en el flujo común. & 24 & T06, T07 & R1, R3, R4 \\
    T09 & \textbf{Entrenar y ajustar la CNN.} Configuración acotada, mejor punto de control y resultados registrados. & 34 & T08 & R1, R3, R4 \\
    T10 & \textbf{Adaptar MobileNet.} Backbone eficiente, cabeza de regresión y estrategia de ajuste definidos. & 24 & T06, T07 & R1, R3, R4 \\
    T11 & \textbf{Entrenar y ajustar MobileNet.} Experimentos comparables, punto de control y resultados registrados. & 34 & T10 & R1, R3, R4 \\
    T12 & \textbf{Adaptar un transformador visual.} Modelo, cabeza de regresión y configuración compatibles con el protocolo. & 26 & T06, T07 & R1, R3, R4 \\
    \bottomrule
  \end{tabularx}
  \caption[Tareas T01--T12]{Tareas T01--T12, con estimación, dependencias y
  recursos. Elaboración propia.}
  \label{tab:planificacion-tareas-1}
\end{table}

La secuencia inicial es
\(\mathrm{T01}\rightarrow\{\mathrm{T03,T04}\}\rightarrow\mathrm{T05}
\rightarrow\mathrm{T06}\rightarrow\mathrm{T07}\). Después, las tres
familias se desarrollan en ramas independientes: T08--T09, T10--T11 y
T12--T13. Esta separación permite avanzar en una familia mientras otra entrena.

\clearpage
\begin{table}[H]
  \centering
  \scriptsize
  \begin{tabularx}{\textwidth}{C{0.75cm}Y C{0.55cm} L{1.55cm} L{1.65cm}}
    \toprule
    \textbf{Cód.} & \textbf{Tarea y resultado de cierre} & \textbf{h} &
    \textbf{Pred.} & \textbf{Recursos} \\
    \midrule
    T13 & \textbf{Entrenar y ajustar el transformador.} Experimentos comparables, punto de control y registro completo. & 36 & T12 & R1, R3, R4 \\
    T14 & \textbf{Aplicar la evaluación común.} Predicciones, métricas por objetivo y análisis de errores incorporados al terminar cada modelo. & 28 & T07; cierre tras T09, T11, T13 & R1, R3, R4 \\
    T15 & \textbf{Medir el coste de inferencia.} Tamaño, latencia y, cuando sea posible, memoria bajo condiciones documentadas. & 24 & T14 & R1, R3, R4 \\
    T16 & \textbf{Convertir y cuantizar el candidato.} Modelo LiteRT válido y comparación antes y después de la conversión. & 24 & T15 & R1, R3, R5 \\
    T17 & \textbf{Validar la ejecución edge.} Inferencia, latencia y limitaciones medidas en el entorno seleccionado. & 20 & T16 & R1, R3, R5 \\
    T18 & \textbf{Sintetizar la comparación.} Tablas de precisión--coste, interpretación y respuesta a la pregunta de investigación. & 20 & T14, T15, T17 & R1, R3, R6 \\
    T19 & \textbf{Analizar sostenibilidad y ética.} Privacidad, sesgo, seguridad, energía y uso indebido. & 12 & T15 & R1, R2, R6 \\
    T20 & \textbf{Redactar y revisar la memoria.} Texto, figuras y trazabilidad actualizados durante el proyecto. & 60 & T01; cierre tras T18, T19 & R1, R2, R6 \\
    T21 & \textbf{Control y seguimiento.} Revisión semanal de horas, avance, riesgos, acuerdos y replanificación. & 24 & T01 & R1, R2, R6 \\
    T22 & \textbf{Cerrar la entrega académica.} Revisión final, portada, anexos, depósito y trámites FIB aplicables. & 10 & T18, T19, T20 & R1, R2, R6 \\
    T23 & \textbf{Preparar la defensa.} Diapositivas, demostración opcional, ensayos y reserva de aula si procede. & 18 & T22 & R1, R2, R3, R6 \\
    \midrule
    & \textbf{Total} & \textbf{540} & & \\
    \bottomrule
  \end{tabularx}
  \caption[Tareas T13--T23]{Tareas T13--T23, con estimación, dependencias y
  recursos. Elaboración propia.}
  \label{tab:planificacion-tareas-2}
\end{table}

Las mediciones se cierran después de evaluar los modelos:
\(\{\mathrm{T09,T11,T13}\}\rightarrow\mathrm{T14}\rightarrow
\mathrm{T15}\rightarrow\mathrm{T16}\rightarrow\mathrm{T17}\).
La conclusión experimental necesita T14, T15 y T17. El cierre técnico reúne
T18, T19 y la versión final de T20; después se realizan T22 y T23.

\begin{table}[H]
  \centering
  \footnotesize
  \begin{tabularx}{\textwidth}{L{1.2cm}L{2.4cm}Y}
    \toprule
    \textbf{Hito} & \textbf{Fecha objetivo} & \textbf{Criterio observable} \\
    \midrule
    M1 & 02/10/2026 & MVP reproducible: datos, particiones, línea base y métricas. \\
    M2 & 11/12/2026 & Tres familias entrenadas y evaluadas con el mismo protocolo. \\
    M3 & 15/01/2027 & Comparación final, validación edge y memoria técnicamente cerrada. \\
    M4 & 29/01/2027 & Depósito y defensa completados; fecha oficial pendiente de Racó. \\
    \bottomrule
  \end{tabularx}
  \caption[Hitos de control]{Hitos y criterios de control de la planificación.
  Elaboración propia.}
  \label{tab:planificacion-hitos}
\end{table}

\thispagestyle{plain}
\clearpage
\begin{landscape}
\section{Diagrama de Gantt}

Cada columna de la figura~\ref{fig:gantt-planificacion} representa una semana.
Las barras largas muestran la concurrencia de documentación y seguimiento, no
una dedicación diaria continua. La semana 20 corresponde al cierre institucional
y a la defensa provisional.

\begin{figure}[H]
  \centering
  \begin{ganttchart}[
    hgrid,
    vgrid,
    x unit=0.77cm,
    y unit chart=0.32cm,
    y unit title=0.42cm,
    title height=1,
    bar height=0.56,
    group height=0.56,
    milestone height=0.50,
    title/.style={fill=deepblue,draw=white},
    title label font=\color{white}\bfseries\tiny,
    bar label font=\tiny,
    group label font=\tiny\bfseries,
    milestone label font=\tiny\bfseries,
    bar/.append style={fill=upcblue!72,draw=upcblue},
    group/.append style={fill=deepblue!82,draw=deepblue},
    milestone/.append style={fill=amber,draw=amber},
    bar label node/.append style={align=right,text width=5.2cm},
    group label node/.append style={align=right,text width=5.2cm},
    milestone label node/.append style={align=right,text width=5.2cm}
  ]{1}{20}
    \gantttitle{Sep.}{3}\gantttitle{Oct.}{4}\gantttitle{Nov.}{4}\gantttitle{Dic.}{5}\gantttitle{Ene.}{4}\\
    \gantttitlelist{1,...,20}{1}\\
    \ganttgroup{Preparación y MVP}{1}{3}\\
    \ganttbar{T01 Alcance y protocolo}{1}{2}\\
    \ganttbar{T02 Revisión bibliográfica}{1}{5}\\
    \ganttbar{T03 Entorno reproducible}{1}{1}\\
    \ganttbar{T04 Adquisición y validación}{1}{2}\\
    \ganttbar{T05 Flujo de datos}{1}{3}\\
    \ganttbar{T06 Particiones y fugas}{3}{3}\\
    \ganttbar{T07 Línea base}{3}{3}\\
    \ganttmilestone{M1 MVP}{3}\\
    \ganttgroup{Modelos y evaluación}{4}{17}\\
    \ganttbar{T08 CNN: implementación}{4}{4}\\
    \ganttbar{T09 CNN: entrenamiento}{5}{6}\\
    \ganttbar{T10 MobileNet: adaptación}{7}{7}\\
    \ganttbar{T11 MobileNet: entrenamiento}{8}{9}\\
    \ganttbar{T12 Transformer: adaptación}{10}{10}\\
    \ganttbar{T13 Transformer: entrenamiento}{11}{12}\\
    \ganttbar{T14 Evaluación común}{6}{13}\\
    \ganttmilestone{M2 Modelos comparados}{13}\\
    \ganttbar{T15 Coste de inferencia}{14}{14}\\
    \ganttbar{T16 Conversión y cuantización}{15}{16}\\
    \ganttbar{T17 Validación edge}{17}{17}\\
    \ganttbar{T18 Síntesis comparativa}{17}{18}\\
    \ganttbar{T19 Sostenibilidad y ética}{15}{18}\\
    \ganttgroup{Gestión y cierre}{1}{20}\\
    \ganttbar[bar/.append style={fill=green!72,draw=green!85!black}]{T20 Memoria}{1}{19}\\
    \ganttbar[bar/.append style={fill=green!72,draw=green!85!black}]{T21 Seguimiento}{1}{19}\\
    \ganttmilestone{M3 Fin técnico}{18}\\
    \ganttbar{T22 Entrega y trámites}{19}{19}\\
    \ganttbar{T23 Defensa}{19}{20}\\
    \ganttmilestone{M4 Fin FIB}{20}
  \end{ganttchart}
  \caption[Diagrama de Gantt]{Diagrama de Gantt de la planificación temporal.
  Las semanas se numeran desde el 14/09/2026; la documentación y el seguimiento
  se ejecutan de forma concurrente. Elaboración propia.}
  \label{fig:gantt-planificacion}
\end{figure}
\end{landscape}

\clearpage
\section{Seguimiento, obstáculos y planes alternativos}

Cada viernes se compararán horas reales y planificadas, tareas terminadas y
riesgos. Una desviación superior al 15\,\% en horas o a una semana en una tarea
de la cadena de cierre activa una revisión con la dirección. Se reestima primero
el trabajo pendiente; las tareas opcionales nunca desplazan el alcance mínimo.

\begin{table}[H]
  \centering
  \scriptsize
  \begin{tabularx}{\textwidth}{L{2.1cm}Y L{2.55cm} L{2.25cm}}
    \toprule
    \textbf{Obstáculo / señal} & \textbf{Tareas alternativas} &
    \textbf{Efecto temporal} & \textbf{Recursos adicionales} \\
    \midrule
    \textbf{R1. Cómputo insuficiente.} Entrenamientos que exceden una iteración. &
    \textbf{B1:} congelar capas, limitar la búsqueda y reducir resolución o lote,
    conservando la prueba común. & Hasta 18 h de margen; no cambia M3. &
    Acelerador externo; coste, si existe, TBD. \\
    \textbf{R2. Datos o almacenamiento.} Descarga incompleta o falta de espacio. &
    \textbf{B2:} continuar con el piloto y la caché mientras se recupera la copia
    completa; verificar integridad antes de reanudar. & Hasta 12 h; T09, T11 y T13
    no se cierran con el piloto. & Almacenamiento temporal R4. \\
    \textbf{R3. Entorno edge incompatible.} T16 no genera un modelo ejecutable. &
    \textbf{B3:} probar operadores alternativos, seleccionar el siguiente modelo
    compatible y medir LiteRT en CPU, documentando la limitación. & Hasta 16 h;
    puede reducir T17, no la comparación principal. & Equipo alternativo o
    préstamo, TBD. \\
    \textbf{R4. Error superior al esperado.} Ninguna familia mejora claramente
    la línea base. & \textbf{B4:} auditar etiquetas, escalas y pérdidas; ejecutar
    un ajuste acotado y, si persiste, analizar el resultado negativo. & Hasta 14 h;
    no altera la fecha final. & R1, R2 y registros R4. \\
    \textbf{R5. Retraso acumulado.} Se agota el margen o fallan dos hitos. &
    \textbf{B5:} eliminar demostrador y extensiones; conservar línea base, CNN,
    MobileNet, transformer, evaluación común y un candidato edge. & Recupera
    hasta 60 h y protege M3/M4. & No requiere recursos nuevos. \\
    \bottomrule
  \end{tabularx}
  \caption[Planes alternativos]{Obstáculos, señales de activación y planes
  alternativos. Estas tareas no forman parte del Gantt base. Elaboración propia.}
  \label{tab:planificacion-riesgos}
\end{table}

El calendario ofrece aproximadamente 60 horas de capacidad no asignada respecto
al ritmo ordinario de 20 semanas. No aumenta el presupuesto de 540 horas: absorbe
festivos, esperas y la activación de B1--B4. Si la desviación supera esa reserva,
se aplica B5. La profundidad, el texto, el demostrador y una búsqueda extensa de
hiperparámetros siguen siendo opcionales.

M3 marca el final técnico: la pregunta de investigación queda respondida con
resultados trazables. M4 marca el final académico: depósito, trámites y defensa.
Si Racó asigna otra fecha, T22 y T23 se desplazarán manteniendo M3 y el trabajo
técnico ya cerrado. Cada hito exige artefactos reproducibles, coherencia entre
memoria y resultados y actualización del registro de decisiones.

\endgroup
